% ================================================================================================ capítulo 11
\chapter{Imágenes astronómicas}\label{cap:imagenes_teoria}

\begin{resumen}
\textbf{En este capítulo:} qué corrigen el bias, el dark y el flat, por qué hay que emparejarlos con cada imagen, cómo se
escala un dark de otra exposición, cómo se combinan los maestros, cómo se modela una estrella para medir su FWHM y qué
hacen las escalas de brillo y ZScale en el visor.
\end{resumen}

\section{Lo que la cámara agrega a la luz}

Una imagen cruda de una cámara astronómica o de una cámara de fotos mezcla la luz del cielo con defectos del instrumento.
Si $L$ es lo que registra cada píxel,
\begin{equation}
L = \underbrace{\text{bias}}_{\text{lectura}} + \underbrace{I_d\,t}_{\text{corriente oscura}}
  + \underbrace{F \cdot S}_{\text{señal del cielo}} ,
\end{equation}
donde $t$ es la exposición, $I_d$ la corriente oscura de ese píxel y $F$ un factor de respuesta (viñeteo, polvo, sensibilidad
de cada píxel) que multiplica la señal $S$. Calibrar es despejar $S$.

\begin{table}[H]
\centering\small
\begin{tabularx}{\linewidth}{@{}l Y Y@{}}
\toprule
Imagen & Cómo se toma & Qué corrige\\
\midrule
\textbf{Bias} & La exposición más corta posible, con la tapa puesta. & El nivel de lectura de la cámara (offset).\\
\textbf{Dark} & Con la tapa puesta, con la \textbf{misma exposición, ganancia (o ISO) y temperatura} que las luces. & El
bias más la corriente oscura y los píxeles calientes. Se resta.\\
\textbf{Flat} & De una superficie iluminada de forma pareja (el cielo del atardecer, una pantalla blanca), con el
\textbf{mismo foco y filtro}. & El viñeteo y las sombras del polvo. Se divide.\\
\textbf{Dark de flat} & Darks con la exposición de los flats. & Lo que se resta a cada flat antes de combinarlos.\\
\bottomrule
\end{tabularx}
\caption{Imágenes de calibración.}
\end{table}

La imagen calibrada es
\begin{equation}
S = \frac{L - \text{dark}}{\text{flat normalizado}} ,
\end{equation}
donde el dark ya incluye el bias. Si no hay dark pero sí bias, se resta solo el bias.

\section{Por qué hay que emparejarlas}

La corriente oscura crece con el tiempo de exposición y con la temperatura del sensor (en muchos sensores se duplica cada
5 a 7~°C), y su escala en cuentas cambia con la ganancia. Un dark de 30~s no sirve para una luz de 120~s. El flat, en
cambio, no depende de la exposición sino del filtro, del foco y de la posición de la cámara: el polvo que se ve en un
filtro no está en el otro. Por eso TurbuLENS empareja cada luz con los darks de su misma exposición, ganancia,
temperatura (con hasta 5~°C de diferencia), binning y tamaño, y con los flats de su mismo filtro, tamaño y binning.

\subsection{El dark escalado}\label{sec:dark_escalado}

Si no hay darks de la exposición de una luz pero sí bias, se puede fabricar uno. La corriente oscura es proporcional al
tiempo, así que, quitando el bias, un dark de $t_\text{dark}$ segundos se convierte en uno de $t_\text{luz}$:
\begin{equation}
\text{dark}(t_\text{luz}) = \left[\text{dark}(t_\text{dark}) - \text{bias}\right]\frac{t_\text{luz}}{t_\text{dark}}
+ \text{bias} .
\end{equation}
TurbuLENS prefiere escalar hacia abajo un dark \textbf{más largo}: al multiplicar por un factor menor que uno, su ruido
también baja; al escalar hacia arriba uno corto, el ruido se amplifica. Sin bias no se puede escalar: se usa el dark
más cercano tal cual, o solo el bias, y el plan lo avisa.

\section{Los maestros}

Cada grupo de imágenes de calibración se combina, píxel por píxel, en una imagen \textbf{maestra}. Hay tres formas:
\begin{itemize}
\item \textbf{Mediana} (la predeterminada). Ignora los valores raros de un píxel, como un rayo cósmico o un satélite que
cruzó en una sola imagen. Para datos normales su ruido es $\sqrt{\pi/2} \approx 1{,}25$ veces el de la media.
\item \textbf{Media con recorte a 3~σ}. En cada píxel descarta los valores a más de 3 desviaciones robustas de la
mediana y promedia el resto: casi tan robusta como la mediana y casi tan silenciosa como la media.
\item \textbf{Media}. El menor ruido si no hay valores raros.
\end{itemize}
El ruido de un maestro de $N$ imágenes es $\sqrt N$ veces menor que el de una sola, así que conviene tomar muchas (20 o más
darks y flats). Con series grandes, TurbuLENS combina por franjas de filas en un archivo temporal, para no tener todas las
imágenes en la memoria.

\textbf{Flats.} Antes de combinarlos, a cada flat se le resta su dark (o el bias) y se divide entre su propia mediana, para
que flats con brillos distintos pesen igual. El maestro se divide otra vez entre su mediana, de modo que valga 1 en
promedio: dividir por él corrige la forma sin cambiar el nivel. Los píxeles con menos del 5\,\% de la mediana (muertos o
tapados) se dejan en 1, para no dividir por casi cero.

\textbf{Flats saturados.} Si al tomar un flat el sensor llega a su valor máximo, todos los píxeles quedan iguales y el
flat sale plano: no corrige nada. El histograma de un buen flat está cerca de la mitad del rango de la cámara. TurbuLENS
cuenta los píxeles en el valor máximo y avisa si son muchos.

\section{El perfil de una estrella}\label{sec:psf}

La imagen de una estrella, en una exposición larga, se parece a una campana: el \textbf{perfil de punto} (PSF). TurbuLENS
lo modela con una gaussiana de dos dimensiones sobre un fondo constante:
\begin{equation}
I(x,y) = B + A\,\exp\!\left[-\frac{(x-x_0)^2}{2\sigma_x^2} - \frac{(y-y_0)^2}{2\sigma_y^2}\right] ,
\end{equation}
con seis parámetros: el fondo $B$, la amplitud $A$, el centro $(x_0, y_0)$ y los anchos $\sigma_x$ y $\sigma_y$. El ajuste
busca los parámetros que hacen mínima la suma de los cuadrados de las diferencias entre el modelo y los píxeles de la región,
con el método de \textbf{Levenberg–Marquardt}: arranca con los momentos de la imagen (centro de luz y anchos) y en cada paso
combina un paso de Gauss–Newton con uno de descenso por el gradiente, más de uno o del otro según si el paso anterior
mejoró el ajuste (sección~\ref{sec:codigo_gauss}).

\begin{figure}[H]
\centering
\begin{tikzpicture}
\begin{axis}[width=0.8\linewidth,height=5.4cm,xlabel={Distancia al centro (en unidades de $\sigma$)},
  ylabel={Brillo / pico},xmin=0,xmax=3.2,ymin=0,ymax=1.08,axis lines=left,grid=major,grid style={grisclaro},
  tick label style={font=\footnotesize},label style={font=\small},/pgf/number format/use comma,
  ytick={0,0.5,0.607,1},yticklabels={0;0,5;0,61;1},extra x ticks={1,1.1774},
  extra x tick labels={},extra x tick style={grid=none}]
\addplot[acento,line width=1.4pt,domain=0:3.2,samples=80] {exp(-x^2/2)};
\addplot[acentoosc,dashed,line width=1pt] coordinates {(1,0) (1,0.6065)};
\addplot[naranja,dashed,line width=1pt] coordinates {(1.1774,0) (1.1774,0.5)};
\node[anchor=west,font=\footnotesize,text=acentoosc] at (axis cs:0.12,0.68) {$\sigma$: 61\,\% del pico};
\node[anchor=west,font=\footnotesize,text=cobre] at (axis cs:1.25,0.5) {FWHM/2 $= 1{,}177\,\sigma$: 50\,\%};
\end{axis}
\end{tikzpicture}
\caption{Perfil de una gaussiana. A un sigma del centro el brillo cae al $e^{-1/2} = 61\,\%$; a medio FWHM, al 50\,\%.}
\end{figure}

El \textbf{ancho a media altura} es la distancia entre los dos puntos donde el brillo cae a la mitad del pico:
\begin{equation}
e^{-\text{FWHM}^2/(8\sigma^2)} = \tfrac12 \quad\Longrightarrow\quad \text{FWHM} = 2\sqrt{2\ln 2}\;\sigma \approx 2{,}355\,\sigma .
\end{equation}
TurbuLENS reporta $\sigma = \sqrt{\sigma_x\sigma_y}$ (el «radio al primer sigma»), el FWHM correspondiente y la
elipticidad $1 - \sigma_\text{menor}/\sigma_\text{mayor}$, que delata problemas de seguimiento o de óptica. Con la escala
de placa, el FWHM en píxeles pasa a segundos de arco, y en una exposición larga y bien enfocada eso es el seeing.

\begin{nota}[La gaussiana es una aproximación.]
El perfil real de una estrella vista a través de la turbulencia tiene alas más extendidas que una gaussiana (se describe
mejor con una función de Moffat). El núcleo, que es lo que define el FWHM, sí lo describe bien la gaussiana, y su ajuste
es más estable con regiones pequeñas. Para medir el seeing con precisión, el DIMM sigue siendo la referencia.
\end{nota}

\section{Cómo se ve una imagen: escalas y límites}

Una imagen astronómica tiene valores entre, digamos, 1\,000 y 60\,000 cuentas, pero lo interesante (el fondo y las
estrellas débiles) ocupa unas pocas decenas de cuentas cerca del fondo. Para verla, el visor hace dos cosas, sin cambiar
los datos:
\begin{enumerate}
\item \textbf{Límites} $[z_1, z_2]$: los valores por debajo de $z_1$ se ven negros y los de encima de $z_2$, blancos.
\item \textbf{Escala}: cómo se reparte el gris entre los dos. Con $t = (v - z_1)/(z_2 - z_1)$ entre 0 y 1:
lineal, $t$; logarítmica, $\log_{10}(1 + 1000\,t)/3$; raíz, $\sqrt t$; asinh, $\operatorname{asinh}(10\,t)/\operatorname{asinh}10$;
e histograma ecualizado, que reparte los grises para que cada tono tenga la misma cantidad de píxeles.
\end{enumerate}

\textbf{ZScale} es el algoritmo de IRAF que usa DS9 para elegir los límites. Toma una muestra de unos mil píxeles, los
ordena y ajusta una recta a los valores ordenados, rechazando iterativamente los que se alejan más de 2,5 desviaciones
(las estrellas brillantes). La pendiente de esa recta, dividida entre un contraste de 0,25, da un rango centrado en la
mediana: muestra bien el fondo y las estrellas débiles, aunque las brillantes salgan saturadas. Como se calcula imagen por
imagen, dos imágenes casi iguales pueden quedar con límites distintos; para compararlas hay que fijar los mismos límites.
